\section{Inference algorithms}

Contractions of sub-networks are the local inference steps.
We review algorithms like AC-3 in that perspective and give with knowledge propagation a generic inference scheme.


%\sect{Deciding Entailment by Local Contractions}\label{subsec:LocalEntailment}

When having a Constraint Satisfaction Problem on a large number of variables, which are densely connected by constraint tensors, direct exploitation of the global entailment criterion in \theref{the:contCriterionLogEntailment} will be infeasible.
An alternative to deciding entailment by global operations is the use of local operations.
Here we interpret a part of the network (for example a single core) as an own knowledge base (with atomic formulas being the roots of the directed subgraph, that is potentially differing with the atoms in the global perspective) and perform entailment with respect to that.


\subsect{Monotonicity of Entailment}

\todo{Remove markov networks for monotonocity statements.
Understand as the proof of the rule $i$.
}

Vanishing local contractions provide sufficient but not necessary criterion to decide entailment, as we show in the next theorem.

\begin{theorem}[Monotonicity of Entailment]
    \label{the:monotonEntailment}
    For any Markov Network $\probtensor^{\graph}$ on the decorated hypergraph $\graph$ and any subgraph $\secgraph$, we have for any formula that $\probtensor^{\graph}\models\exformula$ if $\probtensor^{\secgraph}\models\exformula$.
\end{theorem}

%To show this theorem, we show the following lemma, that whether a contraction of non-negative tensors vanishes, a vanishing contraction of a subset of these tensors is a sufficient criterion.
To prove the theorem, we first establish the following lemma that states if a contraction of non-negative tensors vanishes, the vanishing of a contraction over a subset of these tensors is a sufficient criterion.

\begin{lemma}
    \label{lem:monotocityOfVanishingContractions}
    For any non-negative tensor network $\extnet$ on $\graph$ and $\secedges\subset\edges$ we have the following.
    For $\secgraph=(\secnodes,\secedges)$ with $\secnodes=\cup_{\edge\in\secedges}\edge$ and the tensor network $\tnetof{\secgraph}$ with tensors coinciding on $\secedges$ with those in $\extnet$ we have
    \begin{align*}
        \contraction{\extnet} = 0
    \end{align*}
    if $\contraction{\tnetof{\secgraph}}=0$.
\end{lemma}
\begin{proof}
    Since the tensor network $\tnetof{\secgraph}$ is non-negative, we have whenever $\contraction{\tnetof{\secgraph}}=0$ that
    \begin{align*}
        \contractionof{\tnetof{\secgraph}}{\catvariableof{\secnodes}} = \zerosat{\catvariableof{\secnodes}} \, .
    \end{align*}
    It follows with the commutation of contractions (see \theref{the:splittingContractions} in \charef{cha:messagePassing}), that
    \begin{align*}
        \contractionof{\extnet}{\catvariableof{\nodes}}
        &= \contractionof{
            \{\hypercoreof{\edge} \wcols \edge\in\edges/\secedges\}
            \cup \{\contractionof{\tnetof{\secgraph}}{\catvariableof{\secnodes}}\}
        }{\catvariableof{\nodes}} \\
        &=    \contractionof{
            \{\hypercoreof{\edge} \wcols \edge\in\edges/\secedges\}
            \cup \{\zerosat{\catvariableof{\secnodes}}\}
        }{\catvariableof{\nodes}} \\
        &= 0
    \end{align*}
    Thus, also the contraction of $\extnet$ vanishes in this case.
\end{proof}

\begin{proof}[Proof of \theref{the:monotonEntailment}]
    We use \lemref{lem:monotocityOfVanishingContractions} on the subset $\tnetof{\secgraph}$ of the cores $\extnet$ to the Markov Network $\probof{\graph}$, which itself defines the Markov Network $\probof{\secgraph}$.
    Whenever $\probof{\secgraph}\models\exformula$ for a formula $\exformula$, then we have by \theref{the:contCriterionLogEntailment}
    \begin{align*}
        \contraction{\tnetof{\secgraph} \cup\{\lnot\exformula\}} = 0 \, .
    \end{align*}
    It follows with \lemref{lem:monotocityOfVanishingContractions} that also
    \begin{align*}
        \contraction{\tnetof{\graph} \cup\{\lnot\exformula\}} = 0 \, .
    \end{align*}
    and therefore $\probof{\graph}\models\exformula$.
\end{proof}


\begin{remark}
    To make use of \theref{the:monotonEntailment} we can exploit any entailment criterion.
    However, there is no general statement about entailment possible, when the local entailment does not hold.
    \theref{the:monotonEntailment} therefore just provides a sufficient but not necessary criterion of entailment with respect to $\probtensor^{\graph}$.
\end{remark}

\subsection{Syntactical inference algorithms}

\todo{
    Clauses are tensors, which either shrink or trivialize when contracted with basis vectors.
}

\subsubsection{Forward Chaining in Horn Logics}

\todo{
    When all clauses definite and of length at most $2$, then both the $\onetuple{\nodes}$ and $\zerotuple{\nodes}$ satisfy the constraints.
    Forward chaining as copying length $1$ clauses to other clauses and contract them.
    When terminated all clauses are of lenght at most $2$, and definite if they initially where definite (then any clause of length $1$ is a fact).
}



\subsubsection{Resolution}

\begin{theorem}[Grand resolution theorem]
    When a constraint tensor network of clauses is unsatisfiable, then the empty clause is in the closure under resolution.
\end{theorem}
\begin{proof}
    We give a standard proof \cite{russell_artificial_2021} with the perspective of contractions.
    Having a clause tensor network we add basis vectors to get a elementary tensor network, which is satisfiable when we started with the closure under resolution.
\end{proof}

\subsubsection{Generic inference rules}

More generally, whenever
\begin{align*}
    \bigwedge_{\exformula\in\formulaset} \exformula \models \secexformula
\end{align*}
then the support of the sub contraction of $\formulaset$ leaving the variables of $\secexformula$ open produces a formula entailing $\secexformula$.

